% !TEX root = ../main.tex

% 中英标题：\chapter{中文标题}[英文标题]
\chapter{基于三元组骨架引导的检索增强生成方法}[Retrieval-Augmented Generation Method Guided by Triplet Skeleton]

\section{引言}[Introduction]

大语言模型（Large Language Models, LLMs）凭借其强大的语义理解能力和海量预训练知识，在多种任务中取得了突破性进展。
然而，在需要实时知识与精确推理的知识密集型任务中，LLMs仍然存在固有局限性。其中最主要的问题在于其知识具有静态性，
受限于训练数据的时间截断，容易导致知识过时。

\begin{figure}[htbp]
  \centering
  \includegraphics[width=0.5\textwidth]{TTRAG_intro.png}
  \caption{LLM在复杂推理任务中的典型工作流程}
  \label{fig：TTRAG_intro}
\end{figure}

为了解决这一问题，检索增强生成（Retrieval-Augmented Generation, RAG）被提出作为一种有效方案。RAG通过从外部知识库中
动态检索相关内容并作为上下文输入，从而缓解知识陈旧问题，并显著提升模型在事实性问答任务中的准确性。
然而，标准RAG方法主要依赖语义相似度进行检索，往往会引入大量无关信息，从而降低上下文的信噪比，影响推理效果。

为了解决文本上下文中信噪比较低的问题，研究者开始尝试将知识图谱（Knowledge Graph, KG）引入RAG框架。
知识图谱通常通过开放信息抽取（Open Information Extraction, OpenIE）从非结构化文本中自动抽取关系，
并以三元组形式组织知识。这种结构化表示能够提供清晰的推理框架，使模型能够更快地捕捉核心信息并沿着明确的推理路径进行推断。
然而，现有的基于知识图谱增强的RAG方法（KG-RAG）仍然面临以下关键挑战：

\begin{itemize}
\item \textbf{多跳推理路径的脆弱性}： 在自动构建且存在噪声的知识图谱上进行多跳推理具有较大难度。
错误关系或缺失连接可能导致推理链条中断，从而影响推理的完整性。

\item \textbf{三元组语义稀疏性}： 尽管知识图谱三元组提供了清晰的结构框架，但其语义信息通常较为简洁，
缺乏丰富的上下文细节与解释能力。这种语义稀疏性限制了模型生成复杂和全面答案的能力。

\item \textbf{图结构利用不足}： 大多数KG-RAG方法仅将知识图谱作为检索索引，获取离散的三元组集合，
并将其简单拼接后输入模型，忽略了图中丰富的拓扑结构信息，而这些结构信息对于复杂推理至关重要。
\end{itemize}

如\figref{fig：TTRAG_intro}所示，针对上述问题，本文提出一种基于三元组骨架引导的文本检索增强生成方法
——\textbf{T}riple-Guided \textbf{T}ext \textbf{R}etrieval-\textbf{A}ugmented \textbf{G}eneration（\textbf{TTRAG}）。
与传统RAG方法直接检索非结构化文本，或KG-RAG方法使用离散三元组不同，TTRAG提出了一种“两阶段融合”框架。
该方法首先基于查询构建高质量的三元组推理骨架，然后在此结构化骨架基础上引入精确的文本上下文进行补充。

\begin{figure}[htbp]
  \centering
  \includegraphics[width=0.98 \textwidth]{TTRAG_overview.pdf}
  \caption{TTRAG整体框架}
  \label{fig：TTRAG_overview}
\end{figure}

\figref{fig：TTRAG_overview}展示了模型TTRAG的整体框架，
TTRAG 专为复杂的推理任务而设计。它首先构建知识三元组的推理骨架，以确保逻辑严密的推理路径。
然后，它使用相关的文本内容对该路径进行索引，从而提供丰富的语义信息，实现高效准确的推理。
与其他 RAG 方法相比，TTRAG 具有部分可解释性，并提供清晰、渐进的推理路径。

这种“先构建骨架、再进行语义增强”的策略，不仅为模型提供了清晰且可验证的推理路径，同时弥补了结构化知识语义表达不足的问题，
实现了结构化知识与非结构化文本的有效融合。TTRAG方法的主要优势如下：

\begin{itemize}
\item \textbf{高精度上下文融合}： 通过先构建低噪声的三元组推理骨架，再引入与骨架高度相关的文本信息，
实现结构化知识与非结构化知识的高效融合，从而显著降低无关信息干扰。

\item \textbf{可解释且可靠的推理过程}： 基于显式的三元组推理链，模型能够生成透明且可验证的推理路径，
有效提升生成结果的可信度并减少幻觉现象。

\item \textbf{优越的实验性能}： 在多个多跳问答数据集上的实验结果表明，TTRAG在准确率和F1指标上均优于现有方法，
且具有良好的泛化能力。

\item \textbf{上下文信息密度提升}: 在三个多跳问答数据集上的实验表明，相较于基线方法，TTRAG可减少约70\%--80\%的提示词长度，
显著降低了上下文中的噪声干扰并提升其知识密度。
\end{itemize}

% 待修改
本章组织结构如下：3.2节介绍候选生成-候选排序两阶段实体标准化方法。3.3节介绍基于多层次实体表示的相似度计算方法。
3.4节介绍模型训练方法。3.5节介绍实验结果与分析。3.6节为本章小结。

\section{文档检索与知识图谱构建}[Document Retrieval and Knowledge Graph Construction]

如\figref{fig：TTRAG_framework}所示，TTRAG首先从文档集合中检索出与查询最相关的前 $k$ 个文档，
之后通过上下文学习提示大语言模型作为知识图谱生成器 $g$，对每个文档进行OpenIE，从这些文档中抽取实体关系三元组，
从而构建文档级知识图谱。最后聚合这 $k$ 个文档，形成与查询相关的完整知识图谱 $\mathcal{G}_D$。

\subsection{预备知识}[Preliminary]

\paragraph{知识图谱表示}

知识图谱（Knowledge Graph, KG）可形式化表示为一个有向图 $\mathcal{G} = (\mathcal{E}, \mathcal{R}, \mathcal{T})$。
其中，$\mathcal{E}$ 表示实体集合（节点），$\mathcal{R}$ 表示关系集合（边类型），
$\mathcal{T} \subseteq \mathcal{E} \times \mathcal{R} \times \mathcal{E}$ 表示事实三元组集合。
每一个三元组 $(h, r, t)$ 表示从头实体 $h \in \mathcal{E}$ 到尾实体 $t \in \mathcal{E}$ 的一条有向关系边，
其中 $r \in \mathcal{R}$ 表示实体之间的关系。上述三元组构成了知识图谱中结构化知识的基本单元。

\begin{definition}[证据链]
\label{def:evidence_chain}
在知识图谱 $\mathcal{G}$ 中，\textbf{证据链（Evidence Chain）}是由多个相互连接的三元组组成的有序序列，
用于表示一条完整的推理路径。长度为 $l$ 的证据链 $C$，从起始实体 $e_0$ 到终止实体 $e_l$，可表示为：
\begin{equation}
C(e_0, e_l) = \langle (e_0, r_1, e_1), (e_1, r_2, e_2), \dots, (e_{l-1}, r_l, e_l) \rangle
\end{equation}

其中，$e_i \in \mathcal{E}$（$i = 0, 1, \dots, l$），$r_j \in \mathcal{R}$（$j = 1, 2, \dots, l$）。
在该序列中，相邻两个三元组 $(e_{i-1}, r_i, e_i)$ 与 $(e_i, r_{i+1}, e_{i+1})$ 通过共享实体 $e_i$ 相连接，
从而形成连续的推理路径。该结构能够自然地刻画从 $e_0$ 到 $e_l$ 的多跳推理过程，其中每一个三元组均作为支撑整体推理的证据。
\end{definition}

\paragraph{任务定义}

给定自然语言问题 $q$ 以及知识图谱 $\mathcal{G}$，三元组骨架构建的核心任务是从 $\mathcal{G}$ 中检索与问题相关的证据链集合
 $\mathcal{C}_q = \{C_1, C_2, \dots, C_k\}$。在此基础上，定义生成函数：
\begin{equation}
f(\text{LLMs} \mid q, \mathcal{C}_q) \rightarrow \mathcal{A}_q
\end{equation}

其中，$\mathcal{A}_q$ 表示针对问题 $q$ 生成的答案。通过利用结构化的证据链 $\mathcal{C}_q$，
模型能够在显式推理路径的引导下完成答案生成，从而提升生成结果的可解释性与可靠性。

\subsection{文档检索与知识图谱构建}[Document Retrieval and Knowledge Graph Construction]

给定用户查询 $q$ 和文档集合 $\mathcal{D}$，TTRAG 首先使用检索器 $\mathcal{R}$ 从文档集合 $\mathcal{D}$ 中
检索出与查询最相关的前 $k$ 个文档，记为：
\begin{equation}
\mathcal{D}_k = \mathcal{R}(q, \mathcal{D}, k)
\end{equation}

其中，$\mathcal{D}_k = \{d_1, d_2, \dots, d_k\}$ 表示按相关性排序的前 $k$ 个文档集合。

检索器 $\mathcal{R}$ 可以是基于向量相似度的稠密检索模型，也可以是针对检索增强生成任务优化的检索模型。
具体而言，首先将查询 $q$ 与文档集合 $\mathcal{D}$ 编码为高维向量表示：
\begin{equation}
\mathbf{q} = \text{Enc}(q), \quad \mathbf{d}_i = \text{Enc}(d_i)
\end{equation}

随后，通过计算查询向量与文档向量之间的相似度（如余弦相似度）：
\begin{equation}
\text{sim}(q, d_i) = \frac{\mathbf{q} \cdot \mathbf{d}_i}{\|\mathbf{q}\| \cdot \|\mathbf{d}_i\|}
\end{equation}

并选取得分最高的前 $k$ 个文档构成集合 $\mathcal{D}_k$，作为后续知识图谱构建的文本基础。

在获得相关文档集合 $\mathcal{D}_k$ 后，本文利用大语言模型执行OpenIE。
具体而言，采用基于上下文学习的提示方式，将大语言模型作为知识图谱生成器 $g$，从每个文档中抽取实体关系三元组。
知识图谱构建提示模板如\figref{fig:kg_prompt}所示。

\begin{figure}[htbp]
  \centering
  \begin{minipage}{0.9\textwidth}
  \begin{tcolorbox}[title=知识图谱构建提示模板]
  \small
  给定标题与文本，从文本中抽取所有知识三元组，形式为 
  \textless 标题；关系；尾实体\textgreater，其中标题为给定标题，尾实体为文本中的短语，关系表示标题与尾实体之间的语义关系。

  \vspace{5pt}
  \texttt{In-Context Few-shot}

  \vspace{5pt}
  标题：\{Title\}

  文本：\{Text\}

  知识三元组：
  \end{tcolorbox}
  \end{minipage}
  \caption{知识图谱构建提示模板}
  \label{fig:kg_prompt}
\end{figure}

对于每个文档 $d_i \in \mathcal{D}_k$，生成对应的三元组集合：
\begin{equation}
\mathcal{T}_i = g(d_i)
\end{equation}

其中，$\mathcal{T}_i = \{(h, r, t)\}$ 表示从文档 $d_i$ 中抽取的实体关系三元组集合。

通过对所有文档的三元组集合进行聚合，得到完整的三元组集合：
\begin{equation}
\mathcal{T}_D = \bigcup_{i=1}^{k} \mathcal{T}_i
\end{equation}

进一步地，可构建对应的知识图谱：
\begin{equation}
\mathcal{G}_D = (\mathcal{E}_D, \mathcal{R}_D, \mathcal{T}_D)
\end{equation}

其中，$\mathcal{E}_D$ 和 $\mathcal{R}_D$ 分别表示由三元组集合 $\mathcal{T}_D$ 推导得到的实体集合与关系集合。

在构建知识图谱的过程中，本文进一步为每个三元组建立其与原始文本的映射关系。具体而言，对于每个三元组 $(h, r, t)$，
记录其对应的原始文档、段落以及句子级别的位置索引，从而实现三元组与文本之间的可追溯关联。
该设计使得后续方法（见第3.4节）能够基于三元组骨架高效地回溯并引入相关文本信息，实现结构化知识与非结构化文本的深度融合。

\begin{figure}[htbp]
  \centering
  \includegraphics[width=0.98 \textwidth]{TTRAG_framework.png}
  \caption{检索与初始化：模型根据问题提取最相关的文档，然后从文档中提取图信息并构建相应的图。
  推理路径：对图中的路径进行剪枝以构建推理链，并对相关文档进行索引。使用 Llama-3-8B-Instruct 进行束搜索。
  最后构建模型需要回答的上下文}
  \label{fig：TTRAG_framework}
\end{figure}

\section{三元组骨架构建方法}[Reasoning Skeleton Construction]

如\figref{fig：TTRAG_framework}所示，在构建文档级知识图谱 $\mathcal{G}_D = (\mathcal{E}_D, \mathcal{R}_D, 
\mathcal{T}_D)$ 的基础上，本文进一步提出一种基于三元组路径的推理骨架构建方法，用于支持多跳推理过程。
所谓推理骨架，是由一系列有序连接的知识三元组构成的路径结构，能够刻画从问题到答案的逐步推理过程。
TTRAG通过迭代式路径扩展、基于大语言模型的三元组选择以及束搜索策略，实现了高效且可靠的推理骨架构建。


\subsection{三元组表示与编码}[Triple Representation and Encoding]

给定问题 $q$ 及其对应的三元组集合 $\mathcal{T}_D = \{t_1, t_2, \dots, t_n\}$，其中每个三元组表示为 $t=(h,r,o)$，
本文首先将所有三元组编码为稠密向量表示：
\begin{equation}
\mathbf{E}_T = \{ \mathbf{e}_{t_1}, \dots, \mathbf{e}_{t_n} \}
\end{equation}

该编码过程通过预训练语义编码模型实现，并在整个推理过程中复用，从而避免重复计算，提高整体效率。

\subsection{迭代式路径扩展}[Iterative Path Expansion]

推理骨架通过逐步扩展路径的方式构建。在第 $j$ 跳（hop）时，当前部分路径可表示为：
\begin{equation}
\mathcal{P}^{(j)} = \{ t_{p_1}, \dots, t_{p_j} \}
\end{equation}

为了结合已有推理信息与原始问题，本文将当前路径与问题进行拼接，构造新的查询：
\begin{equation}
q^{(j)} = \text{Concat}(\mathcal{P}^{(j)}, q)
\end{equation}

并将其编码为查询向量 $\mathbf{e}_{q^{(j)}}$。随后，通过计算查询向量与候选三元组向量之间的相似度（采用内积形式）：
\begin{equation}
s_i^{(j)} = \mathbf{e}_{q^{(j)}}^\top \mathbf{e}_{t_i}
\end{equation}

得到每个候选三元组的相关性得分。

为避免路径中出现重复三元组，本文对当前路径中已选取的三元组进行掩码处理。随后选取得分最高的前 $m$ 个三元组构成候选集合：
\begin{equation}
\mathcal{C}^{(j)} = \text{Top-}m(\{s_i^{(j)}\})
\end{equation}

该过程实现了基于语义匹配的候选路径扩展。

\subsection{基于大语言模型的三元组选择}[LLM-guided Triple Selection]

在获得候选三元组集合 $\mathcal{C}^{(j)}$ 后，本文进一步引入大语言模型对候选路径进行语义判断。
具体而言，将当前路径 $\mathcal{P}^{(j)}$ 与候选集合构造成多选提示（multi-choice prompt），
并利用开源大语言模型作为三元组选择器。推理骨架构建提示模板如\figref{fig:skeleton_prompt}所示。

\begin{figure}[htbp]
  \centering
  \begin{minipage}{0.9\textwidth}
  \begin{tcolorbox}[title=推理骨架构建提示模板]
  \small
  从候选知识三元组中选择一个最合适的三元组，用于扩展当前推理路径，使其能够逐步回答给定问题。如果当前路径已经足够回答问题，则输出 STOP。请仅输出选择结果。
  
  \vspace{5pt}
  \texttt{In-Context Few-shot}
  
  \vspace{5pt}
  当前路径中的三元组：
  
  \{Existing Knowledge Triples\}
  
  问题：
  
  \{Question\}
  
  候选三元组：
  
  \{Candidate Knowledge Triples\}
  
  下一步选择：
  \end{tcolorbox}
  \end{minipage}
  \caption{推理骨架构建提示模板}
  \label{fig:skeleton_prompt}
\end{figure}

模型输出每个候选三元组的选择概率，以及一个特殊的终止选项 $\text{STOP}$：
\begin{equation}
P(c \mid \mathcal{P}^{(j)}, q), \quad c \in \mathcal{C}^{(j)} \cup \{\text{STOP}\}
\end{equation}

当模型选择 $\text{STOP}$ 时，表示当前推理路径已经完成。

该机制结合了语义相似度与语言模型的推理能力，使路径扩展更加符合自然语言语义逻辑。

\subsection{基于束搜索的并行路径扩展}[Beam Search-based Expansion]

为同时探索多条潜在推理路径，本文采用束搜索（Beam Search）策略，在每一步最多保留 $B$ 条候选路径。

对于每一条未完成路径，从候选集合中选取概率最高的前 $M$ 个扩展分支，并计算路径的累积得分：
\begin{equation}
\text{Score}(\mathcal{P}^{(j+1)}) = \text{Score}(\mathcal{P}^{(j)}) \cdot P(c \mid \mathcal{P}^{(j)}, q)
\end{equation}

同时，为抑制低质量路径的扩展，本文引入阈值 $\delta$，过滤掉概率低于该阈值的候选分支，从而减少噪声传播。

所有扩展路径按照累积得分进行排序，并保留前 $B$ 条路径进入下一轮扩展。

\subsection{路径终止与骨架生成}[Termination and Skeleton Generation]

当满足以下任一条件时，路径被标记为完成：
\begin{itemize}
\item 模型选择 $\text{STOP}$；
\item 达到预设定的超参数最大跳数 $L$。
\end{itemize}

最终得到的推理骨架集合表示为：
\begin{equation}
\mathcal{S} = \{ (\mathcal{P}_1, s_1), \dots, (\mathcal{P}_B, s_B) \}
\end{equation}

其中，$\mathcal{P}_i$ 表示一条推理路径，$s_i$ 表示其对应的置信得分。

上述推理骨架为后续答案生成提供了结构化、可解释的推理依据，同时显著提升了推理过程的可靠性与可控性。

\section{基于三元组骨架的上下文证据索引}[Contextual Evidence Indexing based on Reasoning Skeleton]
\subsection{文本证据索引}[Textual Evidence Indexing]

在第3.2节的知识图谱构建过程中，本文利用大语言模型作为知识图谱生成器 $g$，
对检索得到的文档集合 $\mathcal{D}_k$ 构建文档级知识图谱。
具体而言，对于每个文档，构建对应的知识图谱子图，并为每个三元组记录其来源信息，包括原始文档、段落以及句子级别的位置索引。

在推理骨架 $\mathcal{S}$ 构建完成后，本文进一步引入文本证据以增强语义表达能力。
对于骨架中的每一个三元组 $\tau = (h, r, t)$，定义其对应的文本映射函数为：
\begin{equation}
\mathcal{M}(\tau) = \{ d_\tau, p_\tau, s_\tau \}
\end{equation}

其中，$d_\tau$ 表示三元组来源文档，$p_\tau$ 表示对应段落，$s_\tau$ 表示包含该三元组语义信息的文本信息。

基于该映射关系，可以从原始文档中精确定位与三元组相关的文本证据。
进一步地，对于一条推理路径 $\mathcal{P} = \{ \tau_1, \tau_2, \dots, \tau_l \}$，其对应的文本证据集合定义为：
\begin{equation}
\mathcal{E}(\mathcal{P}) = \bigcup_{i=1}^{l} \mathcal{M}(\tau_i)
\end{equation}

\figref{句子过滤图.png} 展示了检索到的原始文本块中过滤出仅包含知识的句子的方式。
通过仅保留包含核心知识的句子，本文方法在保证语义完整性的同时，能够显著减少无关文本信息干扰，提高知识密度，
从而大辐减少上下文长度，即最终输入LLM的提示词长度,为后续答案生成提供更加精确且可解释的输入。

\begin{figure}[htbp]
  \centering
  \includegraphics[width=0.98 \textwidth]{句子过滤图.png}
  \caption{利用三元组骨架过滤原始文本段落得到包含知识的句子}
  \label{句子过滤图.png}
\end{figure}

\subsection{证据链编排}[Evidence Chain Organization]

在完成推理骨架构建与上下文证据索引后，TTRAG 将结构化知识与文本证据进行融合，
并与用户问题共同构建提示（Prompt），输入至大语言模型以生成最终答案。该过程可形式化表示为：
\begin{equation}
\mathcal{A}_q = f_{\text{LLM}}(q, \mathcal{S}, \mathcal{C})
\end{equation}

其中，$f_{\text{LLM}}(\cdot)$ 表示作为答案生成器的大语言模型，
$\mathcal{A}_q$ 表示模型输出的最终答案，$q$ 表示用户查询，
$\mathcal{S}$ 表示推理骨架集合，$\mathcal{C}$ 表示基于推理骨架索引得到的上下文证据文本。

与传统基于图的RAG方法不同，后者通常仅对检索到的子图进行线性拼接或简单摘要，忽略了证据链的逻辑结构，
同时引入大量冗余信息。本文提出一种结构引导的协同推理机制，使大语言模型能够在显式推理路径的约束下进行动态、渐进式推理，
从而提升答案的准确性与连贯性。

在获得推理骨架及其对应的文本证据后，本文首先对证据链进行结构化编排。具体而言，按照推理路径中三元组的顺序，
对其对应的文本片段进行组织与拼接，从而构建逻辑连贯的上下文表示：
\begin{equation}
\mathcal{C} = \text{Concat}(s_{\tau_1}, s_{\tau_2}, \dots, s_{\tau_l})
\end{equation}

与传统基于检索文档的RAG方法直接拼接完整检索文档不同，该过程将离散的结构化知识与非结构化文本信息进行融合，从而有效过滤冗余信息，
使推理过程以自然语言形式呈现。该策略能够显著提升上下文的知识密度，降低无关信息的干扰，
同时大幅减少输入到大语言模型中的提示词长度，从而为后续问答提供清晰且可追溯的证据支持。

\subsection{最终答案生成}[Final Answer Generation]

在构建完成提示后，输入大语言模型进行答案生成。该过程包含以下两个阶段：
\begin{itemize}
\item 首先，大语言模型对当前证据进行充分性判断，即评估推理骨架与文本证据是否能够完整回答用户问题。当证据充足时，
模型基于结构化推理路径生成一个全面且符合逻辑的答案。
\item 当模型判断当前证据不足以支持完整推理时，可以触发进一步的推理与探索过程（如扩展推理路径或补充证据）。
在多轮推理仍未获得明确结论的情况下，模型将综合已有证据与其内在知识，生成最合理的答案。
\end{itemize}

上述机制实现了基于结构化证据的动态问答过程，有效提升了答案生成的准确性、连贯性与鲁棒性。

\section{本章小结}[Brief summary]


